\begin{textsample}{NEG Dim 6 – global\_south\_2025\_03 – Score 1.00 – global\_south\_2025\_03/121884531.txt}  \label{ex:f6_neg_146}
As Ramadan begins in Gaza, the holy month, usually a time of spiritual reflection and community, is overshadowed by devastation, loss, and uncertainty.
While a fragile truce between Israel and Hamas remains in place, the people of Gaza continue to face immense hardships, with widespread destruction and an unpredictable ceasefire casting a shadow over their observances.
Advertisement * Scroll to continue Image : AP Image : AP According to reports from Al Jazeera and The Guardian, last year ' s Ramadan was marked by severe hunger and displacement.
Families struggled to break their fasts with whatever little food was available, often sharing a single can of hummus or beans.
While this year ' s truce has brought a temporary reprieve from violence, the region remains deeply scarred by the war.
War ' s Impact on Religious Observances Much of Gaza ' s religious infrastructure has been reduced to rubble, making communal prayers difficult.
Al Jazeera reports that the Great Omari Mosque, one of the region ' s most historic places of @ @ @ @ @ @ @ @ @ @ worshippers without a central gathering space.
Advertisement * Scroll to continue Despite this, Palestinians in Gaza are determined to maintain their traditions.
In Khan Younis, residents have decorated the ruins of their homes with festive Ramadan lights, holding onto symbols of hope amid destruction.
Signs of Life Amid the Ruins Even in the midst of war, glimpses of normalcy are returning.
The Guardian reported that hundreds of Palestinians gathered in Rafah for the first iftar of Ramadan, sharing a meal among the debris of shattered buildings.
Negotiations for a second phase of the ceasefire are ongoing in Cairo, involving mediators from Israel, Qatar, Egypt, and the United States.
While Hamas has not participated directly, its position is being conveyed through Egyptian and Qatari officials.
Image : APA Image : AP Economic conditions, however, remain dire.
France24 notes that while some supermarkets, like Hyper Mall in Nuseirat, have reopened, the majority of essential goods remain inaccessible to those who have lost @ @ @ @ @ @ @ @ @ @ ' s toll on families is immense.
Over 48,000 people have been killed in Gaza, according to Al Jazeera, leaving thousands of families mourning the loss of loved ones.
For Fatima Al-Absi, a resident of Jabalia, Ramadan has become a painful reminder of what she has lost."There ' s no husband, no home, no proper food, and no proper life,"she told The Times of Israel.
Her husband and son-in-law were killed in the war, her home was damaged, and the mosque she once attended has been destroyed.
Yet, even in the face of such suffering, many Gazans continue to hold onto hope."We don't want war,"Al-Absi says."We want peace and safety."The ceasefire agreement ' s first phase, which facilitated the release of hostages and prisoners, has ended, with negotiations for the next phase stalled.
Israeli Prime Minister Benjamin Netanyahu has expressed willingness to extend the truce through Ramadan and Passover, but Hamas @ @ @ @ @ @ @ @ @ @ : AP Image : AP Meanwhile, Israel has announced it will cut off humanitarian aid to Gaza due to Hamas ' s refusal to extend the initial ceasefire phase.
This decision has sparked further concern about worsening conditions for civilians.
Despite the destruction, Palestinians in Gaza continue to rely on their faith and resilience."We can't buy lanterns or decorations like we do every Ramadan,"says Fatima Barbakh from Khan Younis,"but we still hold onto hope."

% matched lemmas: 
\end{textsample}
